\gddchapter{Interview sequence}


Below is the guideline for the semi-open guided interview. 
The interview starts with easy warm-up questions and progresses towards harder questions as it goes along. 

The answers were translated from polish back to english.

\section{Introduction}
Remind the participant that this is not a test, rather a joint exploration.

\section{Icebreakers}

\begin{QandA}
   \item Now that you’ve tried both versions, how are you feeling? Was there anything that immediately surprised you?
         \begin{answered}
            Szymon: I have to say that from the descriptions themselves it would be really hard for me to follow where I should go or when I even can go.
            The places are there but it's not as obvious to me.
            It's also because English is a second language, so you really have to read carefully, because if something is listed, but it says that you cannot enter, you know...
            It's easy to miss.
            Yeah, it's easy to miss.
            Yes, but I really like the idea.
            I imagine if there were a completely AI-generated
            like live generation, right?
            So any action you do, you get a response from the language models.
            This is a really fun experience.
            If you could just move around just using voice.
            Integrate it deeper.


            Gabriel: like color maybe so like if we have a description and like all of the locations are like purple or something you know what a map would be good a map yeah and you like unlock certain parts of the map when you go there or when you investigate those areas yeah yeah that would be really helpful I feel you and when you played it was there anything that like immediately surprised you or not really I mean the idea itself is kind of

            Szymon: for me at least, it's very new.
            So that's a first.
            And I think that's the main thing.
            I only focus about what can I say, what should I say for
            Yeah, I feel you, I feel you.
            So I focus on that.
            \end{answered}

   \item If you had to describe the difference between these two games to a person who hasn't played either, what would you say?
         \begin{answered}
               Szymon: Well, the second with the manual control, so to speak, well, it's way easier.
               I kind of have a short attention span.
               I mean, most of us do right now, but
               reading a whole paragraph is and then trying to figure out what you what commands should you use yeah by voice is way harder than just going through the list of what you can do yeah and just going for everyone would you say it would be a bit different if the descriptions were more factual and less like not poetic but like if it was like you stand in a room there lies a
               played Disco Elysium?

               Gabriel: No I haven't.

               Szymon: Okay but there's a lot of reading there.
               But you don't have to read those because there's a narrator.
               So you have a part with the description written to you by a narrator and then you have text of what is important, right?
               So he reads you like the bigger story and then you have only the distilled important bits.
               I mean, you still have it on text, but you don't have to follow it by eye, you can just listen to it.

         \end{answered}
         
\end{QandA}




\section{Grand tour questions}
(remember to pause here 3-5 seconds)

\begin{QandA}


   \item From the moment you first used your voice to navigate, what were you thinking as you decided what to say to the system?

         \begin{answered}
            Szymon: Will the game understand me?

            Yeah, that was the first question.

            Yeah, sometimes the game can understand you pretty well and sometimes it doesn't understand you for shit.

            That was the priority.

            What I wanted to do in the game was secondary.

            Gabriel: But the way you were framing, you were trying to think like, okay, will it get me?
            
            Szymon: Yeah, because that's my experience with voice-operated stuff.
         \end{answered}

    \item How was it like to use the voice input in the game?
         \begin{answered}
            Szymon: I mean, most of the time I didn't think it got what I wanted to say, only the most obvious commands, right?
            So, enter this location that's listed and you have to say it.
            At least I did call it exactly as it was listed.
            So, enter food court, not enter... Yeah, like we go to the restaurant.

            Gabriel: We had this situation with the corridor, for instance, right?

            [inaudible]
         \end{answered}


\end{QandA}




\section{Core comparative questions}

\begin{QandA}
   \item In which version did you feel more 'inside' the story world?
         \begin{answered}
            I think the second one because it was easier to navigate so I could focus on the actual goal of why I'm there.
            Yeah.
            So I think that was easier to follow the story.
            Yeah.
            Because if you don't know, you don't have to think thrice to enter a different room or go to a bathroom or something like that.
            But I really like the idea itself.
            I'm gonna be honest.
           
         \end{answered}

   \item How did the 'effort' of thinking of what to say naturally compare to the 'effort' of knowing which keys to press?

         \begin{answered}
            Szymon: Oh, it was way harder.

            Gabriel: Way harder to speak?

            Szymon: Yeah, way harder for, as I said, the first worry was that the game won't understand me, so I really had to like triple check to say the most basic command, but also the most clear one, right?
         \end{answered}

    \item Did the freedom to say 'anything' make you feel more in control, or was it more overwhelming than having a fixed list of keys?

         \begin{answered}
            Szymon: If there was a possibility in this game, so I picked up the crowbar and I asked what's the floor made of, because I just thought that maybe I can just rip all of the floor.
            That would be very cool and that would make me feel very much in control of what I'm doing in the game.

            Gabriel: Okay, I hear you.

            Gabriel: So for you, you would say it entirely depends on the degree of implementation of the voice.
            Okay, I see.
            And in this version of the game,
            you felt more in control in which version?

            Szymon: In the second one, because I knew what were the rules exactly, what are the listed options.

            Gabriel: Yeah.

            Szymon: Yeah.

            Szymon: So that's why.
         \end{answered}
    \item How would you compare the overall experience of using the two game versions? Which one would you choose and why?

         \begin{answered}
            Szymon: Well, the second one, the manual one using the keys is way easier, but I really see the potential for the first one.

             And as I said before,

            If the game was finished and I could actually say just, look at that floor, can I rip it off, can I just smash some windows, do that, I really would enjoy that.
            And just doing that with my voice, I really would like to see, I can say just the whole paragraph, just go on a tirade and my character will actually do it.

             Like I want to do 20 different things and order and then it does them.

            Gabriel: I go there, pick up this item and then I want to rescue the brother and after, let's go finish the game.

            Gabriel: And you go and brew some tea and the game just does it.

            Gabriel: Which of those versions would you choose?

            Szymon: I think the first one because it's...
            I'd like to play with the voice commands on my own.
            I'd just like to figure out what I can say and how I can use it.
            I think the first version is more intriguing.
            It's more engaging.

            
         \end{answered}
\end{QandA}






\section{Probing}
In this part of the interview look into the sticky moments noted above and ask follow-up questions probing for more detail.
Include lexical reflection (eg. you used the word "clunky" when explaining navigation. Could you tell me more about what was happening
in that moment?)

The probing was done as follow up questions to the pre-existing questions and in the discussion during the last question rather than separate questions in this test.

% \begin{QandA}
%    \item It looks like you really value creativity when it comes to playing games. Would you agree with that?
%    \begin{answered}
%       I would agree with that because
%       If I played, let's say, a lot of games, then they can become kind of dull.
%       So if something is new for me, or it's like, you know, creative or interesting,
%       I value that because I'm pretty worn out by some normal or usual type of games like the second one.
%       It was way more similar to the games that I played in the past, I would say.
%    \end{answered}
% \end{QandA}


\section{Final reflection}

\begin{QandA}
   \item After reflecting on everything today, is there anything else you’d like to add about using your voice to play this game that we haven't talked about?

         \begin{answered}
            Szymon: Yeah, I've never had an experience like that.

            Gabriel: unless like tabletop RPGs, right?

            Szymon: So, D\&D, something like that.

            So, that's fun because that's exactly what it is.
            You just say what you do and the Dungeon Master tells you if you can do that or what's the consequence of that.
            And if I could do that, like, just here, I guess, in my home, playing a game by myself, that's a whole new world for me.
            I mean, you hook up VR to that and you're a god, right?

            Gabriel: Yeah, yeah, yeah.

            Gabriel: I'm also with you when it comes to the dynamic generation of the world, right?

            That was my original idea for this.

            I wanted to make a game procedurally.

            Szymon: Procedurally or more like AI-generated, right?

            Gabriel: But I also wanted to create a concrete game logic.
            Like, for instance, you can set a fire to your location, right?
            You wanted to set something on fire, right?
            I wanted to create those kind of systems, right?
            and then force the game to essentially do the creative parts based around AI to create this kind of role-playing experience where you would like, yeah, yeah, yeah.
            hear you on that I feel like that's just like a very very cool idea in general and I also like that the main premise of this for me was that I really wanted to boost I wanted to test for me that the thesis the main like the main question I'm asking here with my thesis is that is this actually more immersive or does it have a potential to be more immersive and how different does the player feel when you use the voice right okay the visual thing when it comes like to video games does the

            Szymon: first thing that advances, right?
            So we have VR now.
            I don't know if you've played Half-Life Alyx.

            Gabriel: I haven't, no.

            I mean, you just walk around.

            You don't have to play the game.

            The story is like secondary.

            You go up to a desk and you pick up all of the items.

            You pick up a pen and just write something on the board.

            I mean, it's insane.

            I played it in a game room once.

            I spent like four hours just doing that.

            I didn't leave the first bed.

            There was a jar with a spider, so I gave him some food.

            t was insane.

            But you hook up something like that, you can do a backflip.

            Do a backflip or kick the door in.

            He does that.

            Szymon: I mean, that's ultimate immersion, right?

            Gabriel: Yeah, yeah.

            Gabriel: There is this game I used to play when I was a kid called Scribblenauts Universe or something like that.

            And it's a little bit like that.
            That was like many years before AI was a thing.
            And in the game you could like, it was a 2D game, you would walk around and you had this magic pen and you could like write anything you want and the game would spawn it.
            the game had a massive library of items and things so you could say oh I want to see Barack Obama here and there would be Barack Obama you know like just super super weird stuff and the game rewarded towards like puzzles you have to like you know there's a group of soldiers you have to distract them somehow so you could like draw a tank and like say I want a tank and you would like run over them or you could say I don't know like I want an ice cream booth and they would go and eat ice cream you know like that it's rewarding the creativity you know yeah yeah so
            That was something I was looking at, and also D\&D as you mentioned, I thought that maybe if you can articulate what you say, that it will boost the feeling that I am doing these things, but I noticed that the main difference is that
            In D\&D you speak in first person just to, I want to do something, I want to do something, but in here you still talk to a machine and the machine like lets you do it or doesn't let you do it, therefore you end up talking to like, for permission, like I want to do this, can we do this, can we do that, and that is not the best experience.
            But what's also interesting is that this is number seven, I think, of all my tests, and half the people, roughly half the people are like, this is more immersive, and half the people this is less immersive.
            So it's a very interesting experience to me looking at.
            And then of course if the implementation was to be like deeper then we could argue but even in a current state there's like a split.

            Szymon: I think in the current state it is less immersive.
            mean if it would be like running really smoothly and also I think there should at least
            imensions of what should a command look like, right?
            Yeah.
            Just like a basic stuff, like you have to say that you want to do a certain thing.

            Gabriel: Yeah.

            Szymon: Not just go like completely free paragraphs.

            Gabriel: Yeah.

            Gabriel: I mean, part of me also wanted to see what would you go for.

            Gabriel: OK, sure.

            Szymon: But yeah, it's an experiment.

            Szymon: Yeah.

            Szymon: If the rules for me would be clear or if I had some more time to figure out what I can say and cannot say for the game to understand me, I think the voice command one would be more immersive.

            Gabriel: I mean there's a lot of factors that are immersion breaking, like for me one of the biggest things is that you have to wait for it and it doesn't work sometimes, right?
            GSo that for me is a big immersion breaker.
            Szymon: That's true.
            Szymon: just look at the processing sign and... And you're in the same spot, I didn't understand you, so... Yeah, yeah, you're like, fuck it.

            Gabriel: Repeating for the first...
            There was one player, she was trying to get into the lounge, but her accent was so bad, so she was saying it like five times, and the game just refused to understand.
            Yeah, so that's... That's an immersion breaker.
         \end{answered}
\end{QandA}

\section{Closing}
Thank the participant for his time and tell them again how much value his input and time gives me.






